
\subsection{A fixed-time comparison of squared gradients}
\label{subsec:trigger-stability}

We compare the squared-gradient statistics on paths already coupled under
the prescribed common schedule of \cref{thm:fixed-schedule-high-prob-r-v15}.
The observation time is deterministic and lies inside an ODE interval.
Define
\begin{equation}
\begin{split}
 \mathcal G_t&:=\int\|\nabla_WF(\mu_t,x)\|^2\,\mu_t(dx),\\
 \mathcal G_t^N&:=\frac1N\sum_{i=1}^N
       \|\nabla_WF(\rho_{\mathbf X_t},X_t^i)\|^2.
\end{split}
\label{eq:def-trigger-functional-squared}
\end{equation}

\paragraph{Recall of the comparison quantities.}
Let $E_{\rm pop}\in\mathcal G_{\rm pop}$ and $K_{\rm pop}$ be the
population moment event and its deterministic moment scale from
\cref{lem:population-moment-event}. Conditional on the population fields
$\mathcal G_{\rm pop}$, the variables $Y_t^i$ are i.i.d. with law $\mu_t$;
see \cref{lemma-gp-jump-implies-conditional-iid}.
For $D\ge1$, recall the sample localization event
\[
 E_{T,D}=\left\{\sup_{s\le T}\max_{i\le N}
          (\|X_s^i\|\vee\|Y_s^i\|)\le D\right\}
\]
from \cref{eq:sample-localization-event}.
The common local coefficients in
\cref{eq:conditional-comparison-envelope,eq:common-local-moment-weights}
and the drift comparison are
\[
 \mathrm L_D=\mathrm L_{K_{\rm pop},D},\qquad
 \gamma_m(D)=\mathrm L_Da_m,\qquad
 \Lambda(D)=\sum_{m\ge1}\gamma_m(D)L_m(1+2D^{m-1}).
\]
Here $a_m$ are the moment weights and $L_m$ are the feature derivative
bounds in \cref{assumption-structured-moment-smoothness}. The existing
discrepancy and moment residual are
\[
 r_t=\left(\frac1N\sum_i\|X_t^i-Y_t^i\|^2\right)^{1/2},\qquad
 \Gamma_t^N=\sum_{m\ge1}\gamma_m(D)
 \left\|\frac1N\sum_i\phi_m(Y_t^i)-\int\phi_m\,d\mu_t\right\|;
\]
see \cref{subsec:tracking-metric,lemma-directional-measure-slot-error}.
Let $B$ be the deterministic ODE velocity bound in
\cref{assumption-uniform-concentration-piecewise-lipschitz}.

\begin{lemma}[Fixed-time discrepancy of the squared-gradient statistic]
\label{lem:fixed-time-trigger-squared-gap}
Assume
\cref{assumption-structured-moment-smoothness,assumption-common-seed,assumption-moment-growth,assumption-uniform-concentration-piecewise-lipschitz},
with the prescribed-schedule coupling in
\cref{definition-ode-and-jump,lem:population-first-finite-coupling}.
The reference conditional product law is
$\operatorname{Law}((Y_t^i)_i\mid\mathcal G_{\rm pop})=\mu_t^{\otimes N}$
by \cref{lemma-gp-jump-implies-conditional-iid}.
Conditional moments are bounded on $E_{\rm pop}$ by
\eqref{eq:population-uniform-conditional-moments};
$E_{T,D}$ is the sample event \eqref{eq:sample-localization-event}.
The existing $B$-Lipschitz ODE condition holds almost surely.
Fix a deterministic
$t\in(0,T)\setminus\{t_1,\ldots,t_{\NofJ}\}$ and $\delta\in(0,1)$.
Except on an event whose intersection with $E_{\rm pop}\cap E_{T,D}$
has probability at most $\delta$, one has
\begin{equation}
 |\mathcal G_t^N-\mathcal G_t|
 \le2B\{(\Lambda(D)+\mathrm L_D)r_t+\Gamma_t^N\}
       +B^2\sqrt{\frac{\log(2/\delta)}{2N}}.
 \label{eq:fixed-time-trigger-squared-gap}
\end{equation}
Equivalently,
\[
 \mathbb P\!\left(
 E_{\rm pop}\cap E_{T,D}\cap
 \left\{|\mathcal G_t^N-\mathcal G_t|>
 2B\{(\Lambda(D)+\mathrm L_D)r_t+\Gamma_t^N\}
 +B^2\sqrt{\frac{\log(2/\delta)}{2N}}\right\}\right)\le\delta.
\]
\end{lemma}
\begin{proof}
\noindent\textbf{Step 1: the velocity bound controls the actual gradients.}
We first work on $E_{\rm pop}$, without imposing $E_{T,D}$.
Recall the ODE in \cref{definition-ode-and-jump} and the population
moment continuity calculation in the proof of
\cref{lemma-common-empirical-moment-residual-directional-rt}.
For two nearby times $s,t$ in the same ODE interval, integrate
$D\phi_m$ along the segment joining the two particle positions.
The first-derivative and conditional moment bounds give
\[
 \sum_m a_m\|M_m(\mu_s)-M_m(\mu_t)\|
 \le B|s-t|\sum_m a_mL_m
       \{1+2(K_{\rm pop}m^{1/\alpha})^{m-1}\}
 \longrightarrow0.
\]
The series is finite by the decay estimate
\cref{lem:conditional-local-coefficient-growth}.
Choose a finite $R\ge1$ containing the finitely many continuous particle
paths in a neighborhood of $t$. On $E_{\rm pop}$,
$\mu_s\in\mathcal P_{\alpha,K_{\rm pop}}$ and
$\rho_{\mathbf X_s}\in\mathcal P_R$, so their moment vectors belong to
$\mathcal V_{K_{\rm pop},R}$ by \cref{eq:localized-feasible-moment-region}.
The local moment and spatial estimates
\cref{eq:local-value-gradient-moment-bound,eq:local-gradient-spatial-bound}
then imply
\begin{align*}
 &\|\nabla_WF(\mu_s,Y_s^i)-\nabla_WF(\mu_t,Y_t^i)\|\\
 &\quad\le \mathrm L_{K_{\rm pop},R}
 \left\{\sum_m a_m\|M_m(\mu_s)-M_m(\mu_t)\|
                      +\|Y_s^i-Y_t^i\|\right\}\longrightarrow0.
\end{align*}
For the empirical moments the corresponding estimate is
\[
 \sum_m a_m\|M_m(\rho_{\mathbf X_s})-M_m(\rho_{\mathbf X_t})\|
 \le B|s-t|\sum_m a_mL_m(1+2R^{m-1})\longrightarrow0.
\]
The same local estimates therefore give continuity of
$\nabla_WF(\rho_{\mathbf X_s},X_s^i)$ at $t$.
The radius $R$ is used only to establish this local continuity; it
imposes no additional sample event.

Using the integral form of the reference ODE, for $h\downarrow0$,
\[
 \frac{Y_{t+h}^i-Y_t^i}{h}
 =-\frac1h\int_t^{t+h}\nabla_WF(\mu_s,Y_s^i)\,ds
 \longrightarrow-\nabla_WF(\mu_t,Y_t^i).
\]
Every difference quotient has norm at most $B$, so its limit does too.
Applying the identical argument to $X$ gives, on $E_{\rm pop}$,
\[
 \|\nabla_WF(\mu_t,Y_t^i)\|\le B,\qquad
 \|\nabla_WF(\rho_{\mathbf X_t},X_t^i)\|\le B.
\]
The population bound extends to $\mu_t$-almost every point. Indeed,
conditional-i.i.d. sampling and $E_{\rm pop}\in\mathcal G_{\rm pop}$ give
\begin{align*}
 0
 &=\mathbb E\!\left[\mathbf1_{E_{\rm pop}}
   \mathbf1_{\{\|\nabla_WF(\mu_t,Y_t^1)\|>B\}}\right]\\
 &=\mathbb E\!\left[\mathbf1_{E_{\rm pop}}
   \int\mathbf1_{\{\|\nabla_WF(\mu_t,x)\|>B\}}\,\mu_t(dx)\right].
\end{align*}
The integrand is nonnegative. Its integral is therefore zero for almost
every population realization in $E_{\rm pop}$, proving
$\|\nabla_WF(\mu_t,x)\|\le B$ for $\mu_t$-almost every $x$.

\par\noindent\textbf{Step 2: compare the gradients at paired particles.}
On $E_{\rm pop}\cap E_{T,D}$, insert the reference empirical moments and use
\cref{eq:local-value-gradient-moment-bound,eq:local-gradient-spatial-bound}:
\begin{align*}
 &\|\nabla_WF(\rho_{\mathbf X_t},X_t^i)
                -\nabla_WF(\mu_t,Y_t^i)\|\\
 &\le\sum_m\gamma_m(D)
       \left\|\frac1N\sum_j\phi_m(X_t^j)-\int\phi_m\,d\mu_t\right\|
        +\mathrm L_D\|X_t^i-Y_t^i\|\\
 &\le\sum_m\gamma_m(D)\frac1N\sum_j
         \|\phi_m(X_t^j)-\phi_m(Y_t^j)\|
        +\Gamma_t^N+\mathrm L_D\|X_t^i-Y_t^i\|\\
 &\le\Lambda(D)r_t+\Gamma_t^N+\mathrm L_D\|X_t^i-Y_t^i\|.
\end{align*}
For the last step, the line segment lies in the radius-$D$ ball, so
the first-derivative bound gives
\[
 \begin{aligned}
 \|\phi_m(X_t^j)-\phi_m(Y_t^j)\|
 &\le\int_0^1
 \|D\phi_m(Y_t^j+s(X_t^j-Y_t^j))\|\,\|X_t^j-Y_t^j\|\,ds
 \\
 &\le L_m(1+2D^{m-1})\|X_t^j-Y_t^j\|.
 \end{aligned}
\]
Here $1+D^{m-1}\le1+2D^{m-1}$ allows reuse of $\Lambda(D)$.
Cauchy--Schwarz gives
\[
 \frac1N\sum_j\|X_t^j-Y_t^j\|
 \le\left(\frac1N\sum_j\|X_t^j-Y_t^j\|^2\right)^{1/2}=r_t.
\]
For two vectors of norm at most $B$,
$|\|u\|^2-\|v\|^2|=|\langle u-v,u+v\rangle|\le2B\|u-v\|$.
Apply this inequality to the two actual gradients bounded in Step 1
and average the preceding estimate:
\begin{equation}
 \left|\mathcal G_t^N-\frac1N\sum_i
          \|\nabla_WF(\mu_t,Y_t^i)\|^2\right|
 \le2B\{(\Lambda(D)+\mathrm L_D)r_t+\Gamma_t^N\}.
 \label{eq:trigger-paired-gradient-comparison}
\end{equation}
No bound on a gradient evaluated at an intermediate measure--point pair
is needed.

\par\noindent\textbf{Step 3: concentration of the reference average.}
For almost every population realization in $E_{\rm pop}$, the scalars
$\|\nabla_WF(\mu_t,Y_t^i)\|^2$ are conditionally i.i.d., lie in $[0,B^2]$,
and have conditional mean
\[
 \mathbb E[\|\nabla_WF(\mu_t,Y_t^i)\|^2\mid\mathcal G_{\rm pop}]
 =\int\|\nabla_WF(\mu_t,x)\|^2\,\mu_t(dx)=\mathcal G_t.
\]
For $B>0$, the bounded-differences inequality
\cref{lemma:mcdiarmid} applied to their average has coordinate
sensitivities $B^2/N$. Consequently
\begin{align*}
 &\mathbb P\!\left(
 \left|\frac1N\sum_i\|\nabla_WF(\mu_t,Y_t^i)\|^2-\mathcal G_t\right|
 >B^2\sqrt{\frac{\log(2/\delta)}{2N}}
 \,\middle|\,\mathcal G_{\rm pop}\right)\\
 &\quad\le
 2\exp\left\{-\frac{2N}{B^4}
       \left(B^2\sqrt{\frac{\log(2/\delta)}{2N}}\right)^2\right\}
 =\delta.
\end{align*}
For $B=0$, both statistics are zero on $E_{\rm pop}$ and the result
holds directly. For $B>0$, the triangle inequality and
\eqref{eq:trigger-paired-gradient-comparison} show that failure of
\eqref{eq:fixed-time-trigger-squared-gap} inside $E_{\rm pop}\cap E_{T,D}$
requires failure of the reference-average bound. Thus its probability
is at most
\[
 \mathbb E\!\left[\mathbf1_{E_{\rm pop}}\,
 \mathbb P\!\left(
 \left|\frac1N\sum_i\|\nabla_WF(\mu_t,Y_t^i)\|^2-\mathcal G_t\right|
 >B^2\sqrt{\frac{\log(2/\delta)}{2N}}
 \,\middle|\,\mathcal G_{\rm pop}\right)\right]
 \le\delta\,\mathbb P(E_{\rm pop})\le\delta.
\]
The conditioning throughout is on the population fields, not on
$E_{T,D}$.
\end{proof}

\begin{remark}[A threshold decision at the same fixed time]
Fix a deterministic threshold $\theta\ge0$. Suppose that
\eqref{eq:fixed-time-trigger-squared-gap} holds and
\begin{equation}
 |\mathcal G_t-\theta|>
 2B\{(\Lambda(D)+\mathrm L_D)r_t+\Gamma_t^N\}
       +B^2\sqrt{\frac{\log(2/\delta)}{2N}}.
 \label{eq:trigger-margin-agreement}
\end{equation}
If $\mathcal G_t>\theta$, then
$\mathcal G_t^N\ge\mathcal G_t-|\mathcal G_t^N-\mathcal G_t|>\theta$.
If $\mathcal G_t<\theta$, then
$\mathcal G_t^N\le\mathcal G_t+|\mathcal G_t^N-\mathcal G_t|<\theta$.
The decisions agree at this prescribed observation time. This comparison
does not identify or synchronize adaptive jump times.
\end{remark}
